% 第11回　プレゼン資料作成
\session{11}{2}{12月9日（水）}{プレゼン資料作成}{AIと構成編集}{yes}

\goals{
  \item 「課題、顧客、コンセプト、空間、体験、効果、実現の道筋」の流れで提案を組み立てられる
  \item 1枚1メッセージのスライドを作れる
  \item AIが出したコピー案から選び、自分たちの言葉に直せる
  \item AIの活用と自分たちの思考の反映を説明するスライドを用意できる
}

\fillin{発表時間　\underline{\hspace{12mm}} 分　／　聞き手}

\work[120mm]{ワーク1}{構成案と、各スライドのメッセージ}
\begin{promptbox}[構成を組み立てる]
企業担当者向けに店舗設計を提案する（発表時間）分のプレゼンの構成を作ってください。聞き手は経営企画と店舗運営の担当者です。各スライドについて、「スライドの見出し（1行のメッセージ）」と「載せる内容」を示してください。提案の要約：（改訂版の要約）
\end{promptbox}
\begin{wstabh}{colspec={Q[c,wd=6mm,m] Q[l,wd=17mm,m] X[3,l,m] X[2,l,m] Q[l,wd=14mm,m]},row{2-Z}={ht=10mm}}
  & 流れ & 見出し（1行のメッセージ） & 載せる内容・使う素材 & 担当 \\
  1 & 課題 & & & \\ 2 & 顧客 & & & \\ 3 & コンセプト & & & \\ 4 & 空間 & & & \\
  5 & 体験 & & & \\ 6 & 効果 & & & \\ 7 & 実現の道筋 & & & \\ 8 & AI活用 & & & \\
  9 & & & & \\ 10 & & & & \\
\end{wstabh}

\work[60mm]{ワーク3}{キャッチコピー（20案から3つ、そして1つ）}
\begin{promptbox}[コピーの案を出す]
次のコンセプトを表すキャッチコピーを20案出してください。10字以内のものと20字以内のものを半分ずつ。企業の担当者に向けた落ち着いた表現と、来店客に向けた親しみやすい表現の両方を含めてください。コンセプト：（コンセプト文）
\end{promptbox}
\begin{wstabh}{colspec={X[2,l,m] X[3,l,m]},row{2-Z}={ht=10mm}}
  残した3案 & 選んだ理由・気になる点 \\
  & \\ & \\ & \\
\end{wstabh}
\answerbox[決定したコピー（自分たちの言葉に直した点も）]{18mm}

\work[70mm]{AI活用スライド}{「AIをどう活用し、自分たちの思考をどう反映させたか」の下書き}
\begin{wstab}{colspec={Q[l,wd=34mm,m] X[l,m]},row{1-Z}={ht=15mm}}
  \hc{AIを使った場面} & \\
  \hc{採用したもの} & \\
  \hc{捨てたもの\newline とその理由} & \\
  \hc{自分たちが\newline 判断したこと} & \\
\end{wstab}

\work[50mm]{リハーサル}{相互フィードバック}
\begin{wstabh}{colspec={X[l,m] X[l,m] Q[c,wd=18mm,m]},row{2-Z}={ht=16mm}}
  良かった点 & 改善点 & 所要時間 \\
  & & 分 \\ & & 分 \\
\end{wstabh}

\begin{nexttime}
  \item スライドを完成させる。「AIをどう活用し、自分たちの思考をどう反映させたか」のスライドを必ず含める
  \item 発表時間内に収まるよう、通しで練習する
\end{nexttime}
\aimemo
